\section{Benchmark Design}
\label{sec:benchmark-design}

This section describes the design of \textsc{RealVuln}~v1.0: the taxonomy of target types it recognizes, how the dataset was constructed and labeled, the algorithm used to match scanner findings against ground truth, and the scoring metrics reported.

%% ----------------------------------------------------------------
\subsection{Target Type Taxonomy}
\label{sec:target-types}

We define five \emph{target types} to classify the code under test.
Version~1.0 of \textsc{RealVuln} covers only Type~1.

\medskip
\noindent\textbf{Type 1: Intentionally vulnerable applications.}
Open-source projects built \emph{on purpose} with security flaws for educational or CTF use (e.g., PyGoat, DVPWA, VAmPI).
These projects offer high vulnerability density and diverse CWE coverage, making them ideal for initial calibration.

\medskip
\noindent Types 2--5 (production CVEs, vulnerable libraries, benchmark roll-ups, and academic reproductions) are planned for future versions and defined precisely in the roadmap (Section~\ref{sec:living-benchmark}).

%% ----------------------------------------------------------------
\subsection{Dataset Construction}
\label{sec:dataset-construction}

\subsubsection{Repository Selection.}
We collected 26 intentionally vulnerable Python repositories from GitHub.
Selection criteria were: (1)~the project must be a multi-file application with a meaningful code structure (single-file toy scripts were excluded), (2)~the project must use a recognizable web framework, and (3)~the project should contribute CWE families not already saturated in the dataset.
Each repository is pinned to a specific commit SHA recorded in \texttt{benchmark-manifest.json}, ensuring that all scanners analyze identical source trees.

The 26~repositories span five web frameworks: Flask (15~repos), Django (3), FastAPI (3), aiohttp (1), and Tornado (1), with 3~repos using lightweight or custom HTTP stacks.
Table~\ref{tab:dataset-stats} summarizes the dataset.

\begin{table}[t]
  \centering
  \caption{Dataset statistics for \textsc{RealVuln}~v1.0.}
  \label{tab:dataset-stats}
  \footnotesize
  \begin{tabular}{@{}lr@{}}
    \toprule
    \textbf{Statistic} & \textbf{Value} \\
    \midrule
    Repositories          & 26  \\
    Total findings        & 796 \\
    \quad Vulnerabilities & 676 \\
    \quad FP traps        & 120 \\
    Frameworks            & 5   \\
    CWE families          & 18  \\
    Language              & Python \\
    Target type           & Type~1 \\
    \bottomrule
  \end{tabular}
\end{table}

\subsubsection{Ground Truth Labeling.}
Every finding in the ground truth was produced by manual review.
Each labeled entry records the following fields:

\begin{itemize}
  \item \textbf{Identity:} a unique \texttt{id} and a boolean \texttt{is\_vulnerable} flag.
  \item \textbf{CWE classification:} a \texttt{primary\_cwe} reflecting the most precise weakness, and an \texttt{acceptable\_cwes} list enumerating alternative CWE identifiers that a scanner may reasonably report for the same flaw (e.g., both CWE-89 and CWE-943 for a SQL injection that uses an ORM raw query).
  \item \textbf{Location:} the file path and a \texttt{start\_line}/\texttt{end\_line} range pinpointing the vulnerable code.
  \item \textbf{Severity:} one of \texttt{critical}, \texttt{high}, \texttt{medium}, or \texttt{low}.
  \item \textbf{Evidence:} the source of the annotation (\texttt{manual\_review}, \texttt{cve\_id}, or \texttt{walkthrough}) and a free-text description explaining \emph{why} the code is or is not vulnerable.
\end{itemize}

\noindent
The full ground-truth schema and example entries are available in the repository.\footnote{\url{https://github.com/kolega-ai/RealVulnBenchmark/blob/main/ground-truth/realvuln-djangoat/ground-truth.json}}

\subsubsection{Authorship Metadata.}
All 26~repositories are tagged \texttt{human\_authored} with high confidence.
These projects predate the widespread availability of large language models for code generation, which is relevant for mitigating data contamination concerns when evaluating LLM-based scanners: the ground truth code was not plausibly generated by the models under test.

\subsubsection{False Positive Traps.}
Of the 796~labeled findings, 120 (15.1\%) are \emph{FP~traps}: code patterns that appear suspicious but are demonstrably safe upon manual analysis.
For example, a login function that passes user input to SQLAlchemy's \texttt{filter\_by()} method superficially resembles SQL injection, but the ORM automatically parameterizes the query.
If a scanner flags such an entry, the match is classified as a false positive and penalized in scoring.
FP~traps serve two purposes: (1)~they test a scanner's ability to distinguish truly vulnerable code from safe-but-suspicious patterns, and (2)~they provide true negatives for computing specificity-related metrics such as the false positive rate.

%% ----------------------------------------------------------------
\subsection{Matching Algorithm}
\label{sec:matching}

Each scanner finding is matched against the ground truth on three fields: (1)~\textbf{file path} (exact match after normalization), (2)~\textbf{CWE} (the scanner's CWE must appear in the ground truth entry's \texttt{acceptable\_cwes} list), and (3)~\textbf{line number} (within $\pm$10 lines of the ground truth range). When a finding matches multiple entries, we prefer real vulnerabilities over FP traps, giving the scanner the benefit of the doubt. Each ground truth entry is consumed at most once. Unmatched findings are false positives; unmatched vulnerable entries are false negatives; unmatched traps are true negatives.

\paragraph{Alternative locations.}
Some scanners report attack chains rather than single-point locations: a single audit result may reference multiple files (the vulnerable handler, its route definition, model layers, etc.) while describing one or more vulnerabilities. To handle this fairly, a finding may declare \emph{alternative locations}: additional (file,~line) pairs that the matcher tries if the primary location does not match. If any location matches, the finding is not a false positive. When a single finding's alternatives match multiple distinct ground truth entries, each match is credited as a separate true positive, since the scanner genuinely identified multiple vulnerabilities in one report. This prevents penalizing chain-of-evidence reporting styles while still counting one false positive per unmatched finding.

The full matching implementation is available in the repository.\footnote{\url{https://github.com/kolega-ai/RealVulnBenchmark/blob/main/scorer/matcher.py}}

%% ----------------------------------------------------------------
\subsection{Scoring}
\label{sec:scoring}

\subsubsection{Primary Metric: $F_3$ Score.}

Two base metrics underpin all scoring:

\begin{itemize}
  \item \textbf{Precision:} of everything the scanner flagged, what fraction was actually vulnerable? High precision means less noise for analysts to triage.
  \item \textbf{Recall:} of all real vulnerabilities, what fraction did the scanner find? High recall means fewer missed vulnerabilities.
\end{itemize}

\noindent
In security, a single missed vulnerability can lead to a breach, while false positives cost analyst time but do not directly create risk. We therefore use a recall-weighted primary metric: the $F_3$ score, scaled to $[0, 100]$:
\begin{equation}
  F_3 = \frac{10 \cdot P \cdot R}{9P + R} \times 100
  \label{eq:f3}
\end{equation}
The $F_\beta$ family generalizes the harmonic mean with parameter $\beta$. Setting $\beta = 3$ weights recall nine times more heavily than precision. We also report $F_2$ ($\beta{=}2$, recall weighted $4\times$) for lower-risk contexts. The full leaderboard data in the repository allows re-ranking under any $F_\beta$.

\subsubsection{Secondary Metric: $F_2$ Score.}
We additionally report the $F_2$ score:
\begin{equation}
  F_2 = \frac{5 \cdot P \cdot R}{4P + R} \times 100
  \label{eq:f2}
\end{equation}
Setting $\beta = 2$ weights recall four times more heavily than precision. $F_2$ remains appropriate for lower-risk projects or development-time scanning where noise reduction matters more, but $F_3$ is the recommended default for production and high-risk environments.

\subsubsection{Per-CWE-Family and Per-Severity Breakdowns.}
Metrics are also computed per CWE family (18 families grouping related CWEs, e.g., SQL Injection covers CWE-89, CWE-564, CWE-943) and per severity level (critical, high, medium, low), revealing category-level blind spots.

\subsubsection{Scoring Modes and Multi-Run Aggregation.}
\textsc{RealVuln} supports two aggregation modes for multi-repo evaluation:

\begin{description}
  \item[\texttt{micro}] pools TP, FP, FN, and TN counts across all repositories before computing metrics. Repositories where the scanner produced no output (e.g., due to a timeout or parsing failure) are simply omitted from the pool.
  \item[\texttt{strict\_micro}] is identical except that repositories with no scanner output are treated as if the scanner reported zero findings: every vulnerable GT entry in that repo becomes a false negative. This mode penalizes incomplete runs and is the conservative primary reporting mode in this paper. All headline scores are F3 (strict).
\end{description}

\noindent
When a scanner is run multiple times on the same repository (e.g., with different prompts or configurations), each run is scored independently.
Per-repo scores are then aggregated across runs using micro-averaging over the pooled confusion matrix.

\subsubsection{Reproducibility.}
The file \texttt{benchmark-manifest.json} locks three pieces of information necessary for exact reproduction: (1)~the content hash of the ground truth directory, (2)~the commit SHA for each of the 26~repositories, and (3)~a hash of the default prompt version used for LLM-based scanners.
Any change to the ground truth labels, repository code, or prompt template produces a new manifest, making it straightforward to detect drift between benchmark versions.
